
This study examined whether baseline NeuroSAT generates heterogeneous violation patterns in near-solutions. Using a dual-metric framework, we observed that the architecture produces distance diversity (d/n range = 0.265, exceeding the 0.20 threshold) but not structural diversity (entropy range = 1.145, below the 2.0 threshold). This asymmetry indicates that deterministic LSTM-based message-passing with a single learned initialization explores different solution quality levels while maintaining structurally similar violation patterns.

The finding identifies a specific architectural limitation: deterministic convergence to uniform violation strategies. Addressing this limitation may require stochastic mechanisms, ensemble initialization, or diversity regularization. The dual-metric framework provides a quantitative method for diagnosing whether an architecture generates diverse near-solutions before implementing multi-stage refinement mechanisms.

Future work includes testing stochastic message-passing variants on the full 80,000-instance G4SATBench dataset with multiple random seeds, implementing ensemble initialization strategies, and exploring diversity regularization methods. Cross-domain validation (e.g., theorem proving, code synthesis with formal specifications) would test whether the dual-metric framework generalizes beyond SAT solving.
